\input{../tex-inputs/latex-leseansicht-vorspann}

\section[Arthur Schnitzler, Lili und Arnoldo Cappellini an Gustav Schwarzkopf, 23. 4. 1928]{L04460 Arthur Schnitzler, Lili und Arnoldo Cappellini an Gustav Schwarzkopf, 23. 4. 1928}
\nopagebreak\mylabel{L04460v}
\rehead{ }\normalsize\beginnumbering\briefempfaengerindex{Schwarzkopf, Gustav@\textsc{Schwarzkopf, Gustav}!zzzCappellini, Arnoldo@\emph{von Arnoldo Cappellini}!1928-04-281@{23. 4. 1928}|(be}\briefempfaengerindex{Schwarzkopf, Gustav@\textsc{Schwarzkopf, Gustav}!zzzCappellini, Lili@\emph{von Lili Cappellini}!1928-04-281@{23. 4. 1928}|(be}\briefempfaengerindex{Schwarzkopf, Gustav@\textsc{Schwarzkopf, Gustav}!zzzSchnitzler, Arthur@\emph{von Arthur Schnitzler}!1928-04-281@{23. 4. 1928}|(be}
\toendnotes[C]{\smallbreak\pagebreak[2]}
\correspDesc{Versand  durch Arthur Schnitzler, Lili Cappellini, Arnoldo Cappellini am 23. 4. 1928 in Istanbul
\newline{}Übermittlung  am 23. 4. 1928 in Istanbul
\newline{}Erhalt  durch Gustav Schwarzkopf im Zeitraum [29. 4. 1928 – 3. 5. 1928?] in Wien}\toendnotes[C]{\smallbreak}
\Standort{DLA, A:Schnitzler, HS.1985.1.1897.}
\physDesc{Bildpostkarte, 130 Zeichen
\newline{}Handschrift Arthur Schnitzler: Bleistift, lateinische Kurrent
\newline{}Handschrift Lili Cappellini: Bleistift, lateinische Kurrent
\newline{}Handschrift Arnoldo Cappellini: Bleistift, lateinische Kurrent
\newline{}Versand: Stempel: »\nobreak{}\oindex{Istanbul@\textbf{Istanbul}|pwk}Stambo\textcolor{gray}{ul}, {[}2{]}3. 4. 28\nobreak{}«.  }\toendnotes[C]{\smallbreak}\pstart{}{\pb}Austriche\oindex{Österreich@\textbf{Österreich}|pw}\pend{}\pstart{}Hr Gustav Schwarzkopf\pend{}\pstart{}Wien I\oindex{I., Innere Stadt@\textbf{I., Innere Stadt}|pw}\pend{}\pstart{}Tiefer Graben 17\oindex{Wien@\textbf{Wien}!I., Innere Stadt@\textbf{I., Innere Stadt}!Tiefer Graben 17@\textbf{Tiefer Graben 17}|pw}\pend{}{\bigskip}
\pstart
           \noindent{}{\pb}\textcolor{gray}{\textbf{Constantinople. Mosquée de Sultan Ahmed\oindex{Sultan-Ahmed-Moschee@\textbf{Sultan-Ahmed-Moschee}|pw}}}\pend
           \vspace{1em}
\pstart
           {\pb}Constantinopel\oindex{Istanbul@\textbf{Istanbul}|pw}{ }23/4 28\pend
           \vspace{0.5em}
\pstart
           Viele herzliche Grüße!\pend
           
\pstart
           Ihr{\\[\baselineskip]}\spacefill\mbox{Arth}{\\[\baselineskip]}\label{T_L04460-1v}\edtext{\spacefill\mbox{{[}hs. L. Cappellini:{]} Lili Cappellini}{ }{\\[\baselineskip]}\spacefill\mbox{{[}hs. A. Cappellini:{]} ACappellini}}{\lemma{\textnormal{\emph{Lili … ACappellini}}}\Cendnote{\textnormal{Die Unterschriften von Lili\pwindex{Cappellini, Lili 13.\,9.\,1909 Wien – 26.\,7.\,1928 Venedig@\textsc{Cappellini, Lili} (13.\,9.\,1909 Wien – 26.\,7.\,1928 Venedig)|pwk} und Arnoldo Cappellini\pwindex{Cappellini, Arnoldo 10.\,8.\,1889 Venedig – 8.\,5.\,1954 Rom@\textsc{Cappellini, Arnoldo} (10.\,8.\,1889 Venedig – 8.\,5.\,1954 Rom)|pwk}
                  befinden sich am rechten Rand des Textfeldes um 90° gedreht.}}}\label{T_L04460-1}\pend
           \leftskip=0em{}\selectlanguage{ngerman}\endnumbering\briefempfaengerindex{Schwarzkopf, Gustav@\textsc{Schwarzkopf, Gustav}!zzzCappellini, Arnoldo@\emph{von Arnoldo Cappellini}!1928-04-281@{23. 4. 1928}|)be}\briefempfaengerindex{Schwarzkopf, Gustav@\textsc{Schwarzkopf, Gustav}!zzzCappellini, Lili@\emph{von Lili Cappellini}!1928-04-281@{23. 4. 1928}|)be}\briefempfaengerindex{Schwarzkopf, Gustav@\textsc{Schwarzkopf, Gustav}!zzzSchnitzler, Arthur@\emph{von Arthur Schnitzler}!1928-04-281@{23. 4. 1928}|)be}\mylabel{L04460h}
\begin{anhang}
\end{anhang}\newcommand{\dateiname}{L04460}\newcommand{\titel}{Arthur Schnitzler, Lili und Arnoldo Cappellini an Gustav Schwarzkopf, 23. 4. 1928}\newcommand{\editorInnen}{Herausgegeben von Jahnke, SelmaMüller, Martin Anton}\input{../tex-inputs/latex-leseansicht-abspann}
